\section{Methods: A Hierarchical Taxonomy}
\label{sec:method}

We organize the core methods of event camera vision into three interrelated pillars: \textit{perception} (\cref{sec:perception}), which addresses tasks that extract structured information from event streams; \textit{reconstruction} (\cref{sec:reconstruction}), which recovers visual signals (images, videos, and 3D scenes) from events; and \textit{understanding} (\cref{sec:understanding}), which connects events with high-level semantics through foundation models and language. For each pillar, we present a structured taxonomy, trace the evolution of approaches, and highlight the paradigm shifts driven by large-scale pre-trained models.


% ======================================================================
% 3.1 Perception
% ======================================================================
\subsection{Event Camera Perception}
\label{sec:perception}

Event camera perception encompasses tasks that extract structured information --- object locations, semantic labels, dense geometry, motion fields, and human poses --- from asynchronous event streams. We survey five core tasks, tracing the evolution from early CNN-based approaches to the current landscape of Transformers, state space models, spiking neural networks, and foundation model adaptations.


% ----------------------------------------------------------------------
% 3.1.1 Object Detection
% ----------------------------------------------------------------------
\subsubsection{Object Detection}
\label{sec:perception_det}

Object detection from event cameras is challenging because events encode brightness changes rather than appearance, and the asynchronous nature of the data requires rethinking the frame-centric detection paradigm.

\noindent\circnum{1}~\textbf{CNN and RNN Approaches.} Early work adapted conventional detectors to event representations. Pseudo-label-based detection~\cite{iacono2018dvs} applied supervised learning on DVS data under ego-motion. RED~\cite{perot2020red} demonstrated detection on a 1-megapixel event camera using a feed-forward CNN, establishing the Gen1 and 1Mpx benchmarks. ASTMNet~\cite{li2022astmnet} introduced an asynchronous spatio-temporal memory network for continuous detection. MatrixLSTM~\cite{cannici2020matrixlstm} proposed a differentiable recurrent surface for processing asynchronous events.

\noindent\circnum{2}~\textbf{Transformer-Based Methods.} RVT~\cite{gehrig2023rvt} marked a turning point by introducing a multi-stage recurrent Vision Transformer that achieved 47.2\% mAP on Gen1 while being 6$\times$ faster than prior art. HMNet~\cite{hamaguchi2023hmnet} proposed a hierarchical neural memory network for low-latency processing. GET~\cite{peng2023get} used group event tokens for efficient representation. SODFormer~\cite{zhao2023sodformer} introduced streaming object detection fusing events and frames. SAST~\cite{peng2024sast} (CVPR'24) employed scene-adaptive sparse Transformers. EGSST~\cite{wang2024egsst} combined graph structure with spatiotemporal Transformers at NeurIPS'24.

\noindent\circnum{3}~\textbf{Graph Neural Networks.} AEGNN~\cite{schaefer2022aegnn} processed events as asynchronous spatiotemporal graphs with per-event message passing. DAGr~\cite{gehrig2024dagr}, published in \textit{Nature}, demonstrated that graph-based event processing achieves the latency of a 5,000-fps camera with the bandwidth of a 45-fps camera for automotive detection --- arguably the most impactful result in event-based perception to date.

\noindent\circnum{4}~\textbf{State Space Models.} S5-ViT~\cite{zubic2024s5vit} (CVPR'24 Spotlight) combined S4/S5-style state space layers with Vision Transformers, achieving 33\% faster training than recurrent methods while being robust to input frequency mismatches. SMamba~\cite{li2025smamba} introduced sparse Mamba for event-based detection at AAAI'25.

\noindent\circnum{5}~\textbf{Spiking Neural Networks.} SNNs offer a biologically plausible and energy-efficient alternative. Spiking-YOLO~\cite{kim2020spikingyolo} was the first SNN-based detector. EMS-YOLO~\cite{su2023emsyolo} enabled deep directly-trained SNNs for detection. SpikeYOLO~\cite{li2024spikeyolo} (ECCV'24 Best Paper Candidate) combined integer-valued training with spike-driven inference. EAS-SNN~\cite{wang2024eassnn} introduced end-to-end adaptive sampling. HsVT~\cite{li2025hsvt} proposed hybrid spiking Vision Transformers at ICML'25.

\noindent\circnum{6}~\textbf{Other Paradigms.} MvHeat-DET~\cite{wang2025mvheat} (CVPR'25) introduced heat conduction-based detection with a mixture-of-experts framework and the EvDET200K benchmark. Ev-3DOD~\cite{li2025ev3dod} pushed event cameras into 3D object detection at CVPR'25.

\noindent\circnum{7}~\textbf{Label-Efficient and Multi-Modal Fusion.} LEOD~\cite{wu2024leod} addressed the label scarcity problem through label-efficient training strategies. EOLO~\cite{zhou2024eolo} fused events and frames for robust all-day detection. CAFR~\cite{cao2024cafr} introduced hierarchical feature refinement for event-frame fusion.

A comprehensive catalog of event-based detection methods is provided in \cref{tab:perception_catalog}; quantitative benchmark comparisons are in \cref{sec:appendix}.


% ----------------------------------------------------------------------
% 3.1.2 Semantic Segmentation
% ----------------------------------------------------------------------
\subsubsection{Semantic Segmentation}
\label{sec:perception_seg}

Semantic segmentation assigns a class label to every pixel in the event stream, enabling dense scene understanding for applications such as autonomous driving and robot navigation.

\noindent\circnum{1}~\textbf{Supervised Approaches.} EV-SegNet~\cite{alonso2019evsegnet} pioneered event-based semantic segmentation with a CNN operating on event histograms. ESS~\cite{sun2022ess} transferred segmentation knowledge from labeled still images to unlabeled event streams via unsupervised domain adaptation, a landmark in reducing annotation requirements. HALSIE~\cite{biber2024halsie} simultaneously exploited image and event modalities. EISNet introduced multi-modal fusion for event-image segmentation. ESEG~\cite{li2025eseg} (AAAI'25) incorporated explicit edge-semantic guidance. HybridFusion (AAAI'25) proposed efficient frame-event fusion.

\noindent\circnum{2}~\textbf{Annotation-Efficient Methods.} HPL-ESS~\cite{liu2024hpless} introduced hybrid pseudo-labeling for unsupervised event segmentation. EV-WSSS~\cite{wang2024evwsss} explored weakly supervised segmentation for event cameras at ECCV'24.

\noindent\circnum{3}~\textbf{SNN and SSM Approaches.} SpikingEDN proposed an adaptive spiking encoder-decoder for efficient segmentation. SLTNet used spike-driven lightweight Transformers. MambaSeg~\cite{chen2024mambaseg} harnessed Mamba for image-event semantic segmentation.

\noindent\circnum{4}~\textbf{Open-Vocabulary and Foundation Model Methods.} OpenESS~\cite{kong2024openess} (CVPR'24 Highlight) was the first to achieve open-vocabulary event segmentation by transferring CLIP knowledge through frame-to-event distillation, achieving 53.93\% mIoU on DDD17 \textit{without any event labels}. EventSAM~\cite{chen2024eventsam} adapted the Segment Anything Model for event streams. OVOSE~\cite{bendali2024ovose} used synthetic events with foundation model distillation. SEAL~\cite{li2026seal} (ICLR'26) enabled open-vocabulary event instance segmentation at both instance and part levels. These methods are further discussed in \cref{sec:understanding}.

\noindent\circnum{5}~\textbf{Instance and Panoptic Segmentation.} GMNN~\cite{yang2023gmnn} applied graph mixer neural networks for panoptic segmentation. CTOS introduced continuous-time object segmentation at high temporal resolution.

\noindent\circnum{6}~\textbf{Motion Segmentation.} Events are naturally well-suited for detecting independently moving objects. EV-IMO~\cite{mitrokhin2019evimo} established the first motion segmentation dataset and pipeline. EMSGC~\cite{zhou2021emsgc} used spatio-temporal graph cuts. Un-EvMoSeg~\cite{zhang2024unevmoseg} achieved unsupervised event-based independent motion segmentation at ECCV'24.

\cref{tab:perception_catalog} catalogs all segmentation methods; quantitative comparisons are in \cref{sec:appendix}.


% ----------------------------------------------------------------------
% 3.1.3 Depth and Optical Flow Estimation
% ----------------------------------------------------------------------
\subsubsection{Depth and Optical Flow Estimation}
\label{sec:perception_depth_flow}

Depth estimation and optical flow estimation are closely related dense prediction tasks that recover scene geometry and pixel-level motion, respectively. Event cameras enable both with high temporal resolution and robustness to challenging illumination --- crucial for real-time 3D perception in autonomous systems. We discuss these two complementary tasks together, as many recent methods address both within unified frameworks.

\paragraph{Monocular Depth Estimation.}
E2Depth~\cite{hidalgo2020e2depth} was the first to learn dense monocular depth from events alone using a recurrent architecture. RAMNet~\cite{gehrig2021ramnet} combined events and frames through a recurrent asynchronous multimodal network. EReFormer~\cite{xu2024ereformer} introduced recurrent Transformers for event-based depth. The most recent trend involves foundation model distillation: EventDAM~\cite{xu2025depthanycvent} (ICCV'25) distills proxy depth labels from Depth Anything v2 to train event-specific networks without ground-truth depth annotations --- a paradigm-defining result that exemplifies the ``era of large models'' approach (see also \cref{sec:understanding}).

\paragraph{Stereo Depth Estimation.}
Stereo event cameras enable disparity estimation from event correspondences. DDES~\cite{tulyakov2019ddes} learned event sequence embeddings for dense stereo. SE-CFF~\cite{nam2022secff} focused on future event predictions for stereo matching. DTC-SPADE~\cite{zhang2022dtcspade} used discrete time convolutions for fast stereo. ADES~\cite{cho2023ades} learned adaptive dense event stereo from the image domain. More recently, EMatch~\cite{li2025ematch} (ICCV'25) proposed a unified framework for both optical flow and stereo matching, and ActiveEventNet~\cite{wang2025activeeventnet} (CVPR'25) introduced active stereo to address texture-less regions.

\paragraph{Optical Flow Estimation.}
The continuous temporal sampling of event cameras makes them natural sensors for optical flow estimation, capturing motion at rates far beyond frame-based cameras.

\noindent\circnum{1}~\textbf{Self-Supervised and Unsupervised.} EV-FlowNet~\cite{zhu2018evflownet} pioneered self-supervised optical flow from events using photometric and smoothness losses. Zhu~\etal~\cite{zhu2019unsupervised} extended this to jointly learn flow, depth, and egomotion without ground truth.

\noindent\circnum{2}~\textbf{Dense Supervised.} E-RAFT~\cite{gehrig2021eraft} adapted the RAFT architecture for events, establishing a strong baseline on MVSEC and DSEC benchmarks. DCEIFlow~\cite{wan2022dceiflow} learned dense and continuous optical flow from events and images.

\noindent\circnum{3}~\textbf{Contrast Maximization.} The Secrets of Event-Based Optical Flow~\cite{shiba2022secretsof} provided deep insights into contrast maximization for flow estimation, later extended to depth and ego-motion~\cite{shiba2024cmaxsecrets}.

\noindent\circnum{4}~\textbf{Architecture Innovations.} TMA~\cite{liu2023tma} aggregated temporal motion for improved flow. RPEFlow~\cite{wan2023rpeflow} fused RGB, point cloud, and events for joint optical and scene flow. EEMFlow~\cite{liu2024eemflow} achieved efficient meshflow estimation at CVPR'24. SDformerFlow introduced spatiotemporal swin spikeformers. BAT~\cite{li2025bat} proposed bidirectional adaptive temporal correlation.

\noindent\circnum{5}~\textbf{SNN-Based.} Spike-FlowNet~\cite{lee2020spikeflownet} demonstrated energy-efficient hybrid spiking-analog flow estimation.

\cref{tab:perception_catalog} catalogs all depth and optical flow estimation methods; quantitative comparisons are in \cref{sec:appendix}.


% ----------------------------------------------------------------------
% 3.1.4 Object Tracking
% ----------------------------------------------------------------------
\subsubsection{Object Tracking}
\label{sec:perception_track}

Event cameras enable tracking at speeds and dynamic ranges impossible for frame-based trackers, making them valuable for high-speed and low-light scenarios.

\noindent\circnum{1}~\textbf{Feature Tracking.} EKLT~\cite{gehrig2019eklt} asynchronously tracked features extracted from frames using events, achieving high temporal resolution. Data-driven approaches~\cite{messikommer2023etracker} learned feature tracking for event cameras. BlinkTrack~\cite{wang2024blinktrack} achieved feature tracking over 100 FPS by fusing events and images.

\noindent\circnum{2}~\textbf{Single Object Tracking.} The field has developed rapidly with dedicated benchmarks. VisEvent~\cite{wang2023visevent} established an RGB-event tracking benchmark. STNet~\cite{zhang2022stnet} applied spiking Transformers for event-only tracking. CEUTrack~\cite{tang2022ceutrack} and COESOT unified color-event tracking. HDETrack~\cite{wang2024hdetrack} (CVPR'24) introduced a high-resolution event tracking benchmark with the EventVOT dataset.

\noindent\circnum{3}~\textbf{Mamba-Based Tracking.} A recent trend applies state space models for efficient temporal reasoning. Mamba-FETrack~\cite{wang2024mambafe} introduced frame-event tracking via Mamba. MambaEVT~\cite{wang2024mambaevt} applied Mamba to pure event-stream tracking.

\noindent\circnum{4}~\textbf{Cross-Modal and Long-Term.} CRSOT~\cite{tang2024crsot} addressed cross-resolution tracking with unaligned sensors. AMTTrack and FELT~\cite{wang2024felt} established a long-term event tracking benchmark. GS-EVT~\cite{wang2024gsevt} combined Gaussian Splatting with event cameras for cross-modal tracking.

\cref{tab:perception_catalog} catalogs all tracking methods; quantitative comparisons are in \cref{sec:appendix}.


% ----------------------------------------------------------------------
% 3.1.5 Action Recognition and Pose Estimation
% ----------------------------------------------------------------------
\subsubsection{Action Recognition and Pose Estimation}
\label{sec:perception_action_pose}

Action recognition and pose estimation both exploit the temporal richness of event cameras to understand human behavior. We group them together because they share complementary goals: recognizing \textit{what} a person is doing (action) and estimating \textit{how} they are moving (pose). Many applications --- such as sports analysis, rehabilitation, and human-robot interaction --- require both capabilities jointly.

\paragraph{Action and Gesture Recognition.}
The temporal richness and privacy-preserving nature of events make them attractive for action and gesture recognition.

\noindent\circnum{1}~\textbf{Gesture Recognition.} IBM's DVS128 Gesture system~\cite{amir2017dvs128gesture} demonstrated fully event-based gesture recognition at low power, establishing a widely-used benchmark. Spike Gating Flow~\cite{zhang2022spikegatingflow} introduced hierarchical SNNs for online gesture recognition at NeurIPS'22.

\noindent\circnum{2}~\textbf{Action Recognition.} N-EPIC-Kitchens~\cite{plizzari2022nepickitchens} extended EPIC-Kitchens with synthetic events for egocentric action recognition. HARDVS~\cite{wang2023hardvs} and THU-EACT-50~\cite{gao2023thueact} established large-scale event action benchmarks. DailyDVS-200~\cite{wang2024dailydvs200} (ECCV'24) further scaled to 200 categories with 47 participants.

\noindent\circnum{3}~\textbf{Efficient Architectures.} SpikePoint~\cite{ren2024spikepoint} (ICLR'24) achieved efficient point-based SNN action recognition. SpikMamba combined SNNs with Mamba for action recognition. EventMAE~\cite{chen2025eventmae} applied masked autoencoders for point-wise action recognition.

\noindent\circnum{4}~\textbf{Language-Guided.} ExACT~\cite{tang2024exact} (CVPR'24) pioneered language-guided conceptual reasoning for event-based action recognition, using GPT-4-generated captions and CLIP-based alignment. This cross-modal approach is further discussed in \cref{sec:understanding}.

\paragraph{Human Pose Estimation.}
Event cameras capture fast human motions without blur, enabling high-frequency pose estimation critical for sports analysis, rehabilitation, and human-robot interaction.

\noindent\circnum{1}~\textbf{Body Pose.} DHP19~\cite{calabrese2019dhp19} established the first multi-view DVS dataset for 3D human pose estimation. EventCap~\cite{xu2020eventcap} achieved monocular 3D capture of high-speed human motions. EventHPE~\cite{zou2021eventhpe} estimated 3D human pose and shape from events. EventPointPose~\cite{wang2024eventpointpose} explored 3D event representations for pose estimation.

\noindent\circnum{2}~\textbf{Hand Pose.} EventHands~\cite{rudnev2021eventhands} demonstrated real-time neural 3D hand pose estimation from event streams. EvRGBHand~\cite{jiang2024evrgbhand} (CVPR'24) complemented events with RGB frames for hand mesh reconstruction. EventEgoHands~\cite{li2025eventegohands} addressed egocentric event-based 3D hand reconstruction.

\noindent\circnum{3}~\textbf{Egocentric Pose.} EventEgo3D~\cite{millerdurai2024eventego3d} (CVPR'24) captured 3D human motion from egocentric event streams using a head-mounted event camera, later extended to EventEgo3D++~\cite{millerdurai2025eventego3dpp} with improved accuracy and robustness.

\cref{tab:perception_catalog} catalogs all action recognition and pose estimation methods; quantitative comparisons are in \cref{sec:appendix}.

% ============================================================
%  Catalog Tables for the Event Camera Survey
%  This file is \input{} into 3_method.tex
%  Contains: tab:perception_catalog, tab:recon_catalog, tab:understanding_catalog
% ============================================================


% ============================================================
% TABLE 1: Perception Catalog
% ============================================================
\begin{table*}[p]
\centering
\caption{Comprehensive catalog of event camera perception methods. Methods are grouped by task and listed in chronological order within each group. The column \textbf{In} denotes the input modality, \textbf{Out} the prediction target, \textbf{Repr.} the event representation, \textbf{Arch.} the backbone architecture, and \textbf{Cat.} the numbered category used in the narrative.}
\label{tab:perception_catalog}
\resizebox{\textwidth}{!}{
\begin{tabular}{rl l c c c c l c}
\toprule
\textbf{\#} & \textbf{Method} & \textbf{Venue} & \textbf{In} & \textbf{Out} & \textbf{Repr.} & \textbf{Arch.} & \textbf{Task} & \textbf{Cat.} \\
\midrule
\multicolumn{9}{l}{\cellcolor{gray!8}\textit{2.1 Object Detection}} \\
\midrule
%% --- CNN, RNN & Graph Approaches ---
1 & DVS-Detection~\cite{iacono2018dvsdetection} & CVPRW'18 & \iconev & Boxes & EF & CNN & Det. & \circnum{1} \\
\rowalt
2 & YOLE~\cite{cannici2019yole} & CVPRW'19 & \iconev & Boxes & EF & CNN & Det. & \circnum{1} \\
3 & MatrixLSTM~\cite{perot2020red} & NeurIPS'20 & \iconev & Boxes & Learned & RNN & Det. & \circnum{1} \\
\rowalt
4 & ASTMNet~\cite{li2022astmnet} & TIP'22 & \iconev & Boxes & VG & RNN & Det. & \circnum{1} \\
5 & AEGNN~\cite{schaefer2022aegnn} & CVPR'22 & \iconev & Boxes & Graph & GNN & Det. & \circnum{1} \\
\rowalt
6 & RENet~\cite{zhou2023renet} & ICRA'23 & \iconev\iconcam & Boxes & VG & CNN & Det. & \circnum{1} \\
7 & DAGR~\cite{gehrig2024dagr} & Nature'24 & \iconev\iconcam & Boxes & Graph & GNN & Det. & \circnum{1} \\
\rowalt
8 & CDE-Net & ICRA'24 & \iconev\iconcam & Boxes & VG & CNN & Det. & \circnum{1} \\
%% --- Transformer & Attention-Based ---
9 & RVT~\cite{gehrig2023rvt} & CVPR'23 & \iconev & Boxes & VG & ViT & Det. & \circnum{2} \\
\rowalt
10 & SODFormer~\cite{huang2023sodformer} & TPAMI'23 & \iconev\iconcam & Boxes & VG & ViT & Det. & \circnum{2} \\
11 & GET~\cite{peng2024get} & ICCV'23 & \iconev & Boxes & Graph & ViT & Det. & \circnum{2} \\
\rowalt
12 & LEOD~\cite{wu2024leod} & CVPR'24 & \iconev & Boxes & VG & ViT & Det. & \circnum{2} \\
13 & State-Space~\cite{zubic2024s5vit} & CVPR'24 & \iconev & Boxes & VG & SSM & Det. & \circnum{2} \\
\rowalt
14 & SAST~\cite{peng2024sast} & CVPR'24 & \iconev & Boxes & VG & ViT & Det. & \circnum{2} \\
15 & EMF & arXiv'25 & \iconev & Boxes & VG & ViT & Det. & \circnum{2} \\
%% --- SNN & Efficient Detection ---
\rowalt
16 & SFOD & CVPR'24 & \iconev\iconcam & Boxes & Spike & SNN & Det. & \circnum{3} \\
17 & EAS-SNN & ECCV'24 & \iconev & Boxes & Spike & SNN & Det. & \circnum{3} \\
\rowalt
18 & MoE-HeatDet~\cite{zhang2025mvheatdet} & CVPR'25 & \iconev & Boxes & EF & MoE & Det. & \circnum{3} \\
19 & EventFly~\cite{kong2025eventfly} & CVPR'25 & \iconev\iconcam & Boxes & VG & Hybrid & Det. & \circnum{3} \\
\rowalt
20 & OpenESS-Det~\cite{zhang2025ovdetection} & arXiv'25 & \iconev & Boxes & VG & ViT & Det. & \circnum{3} \\
\midrule
\multicolumn{9}{l}{\cellcolor{gray!8}\textit{2.2 Semantic Segmentation}} \\
\midrule
%% --- Supervised & Multi-Modal Fusion ---
21 & EV-SegNet~\cite{alonso2019evsegnet} & CVPRW'19 & \iconev & Masks & EF & CNN & Seg. & \circnum{4} \\
\rowalt
22 & EvDistill~\cite{wang2021evdistill} & CVPR'21 & \iconev\iconcam & Masks & VG & CNN & Seg. & \circnum{4} \\
23 & ESS~\cite{sun2022essseg} & ECCV'22 & \iconev\iconcam & Masks & VG & CNN & Seg. & \circnum{4} \\
\rowalt
24 & EV-Transfer & RA-L'22 & \iconev\iconcam & Masks & EF & CNN & Seg. & \circnum{4} \\
25 & CMX & TITS'23 & \iconev\iconcam & Masks & EF & CNN & Seg. & \circnum{4} \\
\rowalt
26 & HMNet~\cite{hamaguchi2023hmnet} & CVPR'23 & \iconev & Masks & VG & RNN & Seg. & \circnum{4} \\
27 & CMDA & ICCV'23 & \iconev\iconcam & Masks & VG & CNN & Seg. & \circnum{4} \\
\rowalt
28 & HALSIE~\cite{biber2024halsie} & WACV'24 & \iconev\iconcam & Masks & VG & Hybrid & Seg. & \circnum{4} \\
29 & EventSAM~\cite{chen2024eventsam} & CVPR'24 & \iconev & Masks & EF & ViT & Seg. & \circnum{4} \\
\rowalt
30 & SAM-E-Adapter~\cite{li2024sameventadapter} & ICRA'24 & \iconev\iconcam & Masks & EF & ViT & Seg. & \circnum{4} \\
31 & OpenESS~\cite{kong2024openess} & CVPR'24 & \iconev\iconcam & Masks & VG & ViT & Seg. & \circnum{4} \\
\rowalt
32 & MUSES & ECCV'24 & \iconev\iconcam & Masks & VG & CNN & Seg. & \circnum{4} \\
33 & EvSAM~\cite{li2025evsam} & TOMM'25 & \iconev\iconcam & Masks & EF & ViT & Seg. & \circnum{4} \\
\rowalt
34 & SLTNet & IROS'25 & \iconev & Masks & Spike & SNN & Seg. & \circnum{4} \\
%% --- Annotation-Efficient & Adaptation ---
35 & HPL-ESS & CVPR'24 & \iconev & Masks & VG & CNN & Seg. & \circnum{5} \\
\rowalt
36 & EV-WSSS & ECCV'24 & \iconev & Masks & VG & CNN & Seg. & \circnum{5} \\
37 & EvLowLightSeg & WACV'25 & \iconev\iconcam & Masks & VG & CNN & Seg. & \circnum{5} \\
%% --- Motion & Instance Segmentation ---
\rowalt
38 & EV-IMO & IROS'19 & \iconev & Masks & EF & CNN & Seg. & \circnum{6} \\
39 & EMSGC~\cite{zhou2021emsgc} & TNNLS'21 & \iconev & Masks & Graph & -- & Seg. & \circnum{6} \\
\rowalt
40 & OVOSE~\cite{kim2024ovose} & arXiv'24 & \iconev & Masks & VG & ViT & Seg. & \circnum{6} \\
41 & Un-EvMoSeg & ECCV'24 & \iconev & Masks & VG & CNN & Seg. & \circnum{6} \\
\rowalt
42 & OutOfRoom & 3DV'24 & \iconev & Masks & VG & CNN & Seg. & \circnum{6} \\
43 & MouseSIS & ECCVW'24 & \iconev\iconcam & Masks & EF & CNN & Seg. & \circnum{6} \\
\rowalt
44 & SEAL~\cite{li2026seal} & ICLR'26 & \iconev\icontext & Masks & VG & ViT & Seg. & \circnum{6} \\
\midrule
\multicolumn{9}{l}{\cellcolor{gray!8}\textit{2.3 Depth \& Optical Flow Estimation}} \\
\midrule
%% --- Monocular Depth Estimation ---
45 & EMDE & ICRAS'22 & \iconev & Depth & VG & CNN & Mono. & \circnum{7} \\
\rowalt
46 & E2Depth~\cite{hidalgo2020e2depth} & 3DV'20 & \iconev & Depth & VG & CNN & Mono. & \circnum{7} \\
47 & RAM-Net~\cite{gehrig2021ramnet} & RA-L'21 & \iconev\iconcam & Depth & VG & RNN & Mono. & \circnum{7} \\
\rowalt
48 & HMNet-Depth~\cite{hamaguchi2023hmnet} & CVPR'23 & \iconev & Depth & VG & RNN & Mono. & \circnum{7} \\
49 & EReFormer~\cite{cho2024ereformer} & ECCV'24 & \iconev & Depth & VG & ViT & Mono. & \circnum{7} \\
\rowalt
50 & EMoDE~\cite{rao2024emodepth} & arXiv'24 & \iconev & Depth & VG & CNN & Mono. & \circnum{7} \\
51 & OnDevice-Depth & CVPR'25 & \iconev & Depth & VG & CNN & Mono. & \circnum{7} \\
\rowalt
52 & Depth AnyEvent~\cite{wu2025eventdam} & ICCV'25 & \iconev & Depth & VG & Hybrid & Mono. & \circnum{7} \\
53 & FUSE & IROS'25 & \iconev\iconcam & Depth & VG & CNN & Mono. & \circnum{7} \\
\rowalt
54 & DERD-Net & NeurIPS'25 & \iconev & Depth & Raw & CNN & Mono. & \circnum{7} \\
%% --- Stereo Depth Estimation ---
55 & MVSEC-Stereo & RA-L'18 & \iconev & Depth & Learned & CNN & Stereo & \circnum{8} \\
\rowalt
56 & StereoSpike~\cite{ranccon2022stereospike} & arXiv'21 & \iconev & Depth & Spike & SNN & Stereo & \circnum{8} \\
57 & DSGN-Event~\cite{zhang2022dtcspade} & CVPR'22 & \iconev & Depth & VG & CNN & Stereo & \circnum{8} \\
\rowalt
58 & SDE-Net~\cite{cho2022scsnet} & ECCV'22 & \iconev\iconcam & Depth & TS & CNN & Stereo & \circnum{8} \\
59 & EvAC3D & ICCV'23 & \iconev\iconcam & Depth & VG & CNN & Stereo & \circnum{8} \\
\rowalt
60 & E-Stereo-FUDA & ECCV'24 & \iconev & Depth & VG & CNN & Stereo & \circnum{8} \\
61 & ESVO2~\cite{zhou2025esvo2} & TRO'25 & \iconev & Depth & TS & Hybrid & Stereo & \circnum{8} \\
\rowalt
62 & URNet & VI'25 & \iconev & Depth & VG & CNN & Stereo & \circnum{8} \\
%% --- Optical Flow Estimation ---
63 & EV-FlowNet~\cite{zhu2018evflownet} & RSS'18 & \iconev & Flow & EF & CNN & Flow & \circnum{9} \\
\rowalt
64 & Spike-FlowNet~\cite{lee2020spikeflownet} & ECCV'20 & \iconev & Flow & Spike & SNN & Flow & \circnum{9} \\
65 & E-RAFT~\cite{gehrig2021eraft} & 3DV'22 & \iconev & Flow & VG & CNN & Flow & \circnum{9} \\
\rowalt
66 & DCEIFlow~\cite{wan2022dceiflow} & TIP'22 & \iconev\iconcam & Flow & VG & CNN & Flow & \circnum{9} \\
67 & BlinkFlow~\cite{luo2023blinkflow} & IROS'23 & \iconev & Flow & VG & CNN & Flow & \circnum{9} \\
\rowalt
68 & EGMBA~\cite{nunes2025egomotion} & CVPR'25 & \iconev & Flow & Raw & Hybrid & Flow & \circnum{9} \\
69 & MeshFlow & TPAMI'25 & \iconev & Flow & VG & CNN & Flow & \circnum{9} \\
\rowalt
70 & NMGNet & ICRA'25 & \iconev & Flow & VG & CNN & Flow & \circnum{9} \\
\midrule
\multicolumn{9}{l}{\cellcolor{gray!8}\textit{2.4 Object Tracking}} \\
\midrule
%% --- Frame-Event Tracking ---
71 & FE108-Track~\cite{zhang2021fe108} & ICCV'21 & \iconev\iconcam & Tracks & EF & CNN & Track. & \circnum{10} \\
\rowalt
72 & CEUTrack~\cite{tang2022ceutrack} & PR'25 & \iconev\iconcam & Tracks & VG & ViT & Track. & \circnum{10} \\
73 & AFNet~\cite{zhang2023afnet} & CVPR'24 & \iconev\iconcam & Tracks & VG & CNN & Track. & \circnum{10} \\
\rowalt
74 & EventVOT-Track~\cite{wang2024hdetrack} & CVPR'24 & \iconev & Tracks & VG & ViT & Track. & \circnum{10} \\
75 & Asyn-MOT & ICRA'25 & \iconev\iconcam & Tracks & VG & CNN & Track. & \circnum{10} \\
%% --- Event-Only Tracking ---
\rowalt
76 & STNet~\cite{zhang2022stnet} & CVPR'22 & \iconev & Tracks & Spike & SNN & Track. & \circnum{11} \\
77 & HDETrack~\cite{tang2024crsottrack} & TMM'25 & \iconev\iconcam & Tracks & VG & ViT & Track. & \circnum{11} \\
\rowalt
78 & eMoE-Tracker & RA-L'24 & \iconev & Tracks & VG & SSM & Track. & \circnum{11} \\
79 & MambaEVT~\cite{liu2024mambafetrack} & arXiv'24 & \iconev\iconcam & Tracks & VG & SSM & Track. & \circnum{11} \\
\rowalt
80 & HDETrack-V2 & arXiv'25 & \iconev & Tracks & VG & ViT & Track. & \circnum{11} \\
%% --- Long-Term & Cross-Modal Tracking ---
81 & TETrack~\cite{wang2024amttrack} & arXiv'24 & \iconev & Tracks & VG & ViT & Track. & \circnum{12} \\
\rowalt
82 & SiamEFT & Front.~Neuro.'24 & \iconev\iconcam & Tracks & VG & CNN & Track. & \circnum{12} \\
83 & EventFly-Track~\cite{kong2025eventfly} & CVPR'25 & \iconev\iconcam & Tracks & VG & Hybrid & Track. & \circnum{12} \\
\rowalt
84 & TUMTraf-EMOT & arXiv'25 & \iconev & Tracks & VG & CNN & Track. & \circnum{12} \\
\midrule
\multicolumn{9}{l}{\cellcolor{gray!8}\textit{2.5 Action Recognition \& Pose Estimation}} \\
\midrule
%% --- Action & Gesture Recognition ---
85 & DVS-Gesture~\cite{amir2017dvsgesture} & CVPR'17 & \iconev & Labels & EF & CNN & Act. & \circnum{13} \\
\rowalt
86 & PointNet-EVS & CVPRW'19 & \iconev & Labels & Raw & CNN & Act. & \circnum{13} \\
87 & E2(GO)MOTION~\cite{plizzari2022e2gomotion} & CVPR'22 & \iconev\iconcam & Labels & VG & CNN & Act. & \circnum{13} \\
\rowalt
88 & EventTransAct & IROS'23 & \iconev & Labels & VG & ViT & Act. & \circnum{13} \\
89 & Bullying10K & NeurIPS'23 & \iconev & Labels & VG & CNN & Act. & \circnum{13} \\
\rowalt
90 & ExACT~\cite{zhou2024exact} & CVPR'24 & \iconev\icontext & Labels & VG & ViT & Act. & \circnum{13} \\
91 & SpikePoint & ICLR'24 & \iconev & Labels & Raw & SNN & Act. & \circnum{13} \\
\rowalt
92 & HARDVS-Method~\cite{wang2024hardvsmethod} & AAAI'24 & \iconev & Labels & VG & CNN & Act. & \circnum{13} \\
93 & DailyDVS-Method~\cite{wang2024dailydvs200method} & ECCV'24 & \iconev & Labels & VG & CNN & Act. & \circnum{13} \\
\rowalt
94 & PASS & NeurIPS'25 & \iconev & Labels & VG & SSM & Act. & \circnum{13} \\
95 & EventFly-Action~\cite{kong2025eventfly} & CVPR'25 & \iconev\iconcam & Labels & VG & Hybrid & Act. & \circnum{13} \\
%% --- Body Pose Estimation ---
\rowalt
96 & DHP19-Method~\cite{calabrese2019dhp19method} & CVPRW'19 & \iconev & Pose & EF & CNN & Pose & \circnum{14} \\
97 & EventCap~\cite{xu2020eventcap} & CVPR'20 & \iconev\iconcam & Pose & EF & Hybrid & Pose & \circnum{14} \\
\rowalt
98 & EventHPE~\cite{zou2021eventhpe} & CVPR'21 & \iconev & Pose & VG & CNN & Pose & \circnum{14} \\
99 & SSTFormer & TCSVT'23 & \iconev & Pose & Spike & SNN & Pose & \circnum{14} \\
\rowalt
100 & EventEgo3D~\cite{millerdurai2024eventego3d} & CVPR'24 & \iconev & Pose & VG & CNN & Pose & \circnum{14} \\
101 & EHPT-XC~\cite{xu2024ehptxc} & NeurIPS'24 & \iconev\iconcam & Pose & VG & Hybrid & Pose & \circnum{14} \\
\rowalt
102 & E-POSE & Sci.~Data'25 & \iconev & Pose & EF & CNN & Pose & \circnum{14} \\
%% --- Hand & Egocentric Pose Estimation ---
103 & EventHandPose & ICCV'23 & \iconev\iconcam & Pose & EF & CNN & Hand & \circnum{15} \\
\rowalt
104 & EvHandPose & TPAMI'24 & \iconev & Pose & EF & CNN & Hand & \circnum{15} \\
105 & EvRGBHand~\cite{zhou2024evrgbhand} & CVPR'24 & \iconev\iconcam & Pose & VG & CNN & Hand & \circnum{15} \\
\rowalt
106 & InterHand-Event & 3DV'24 & \iconev & Pose & EF & CNN & Hand & \circnum{15} \\
\bottomrule
\end{tabular}
}

\vspace{2pt}
\begin{tcolorbox}[colback=gray!5, colframe=gray!50, boxrule=0.4pt, arc=2pt, left=4pt, right=4pt, top=2pt, bottom=2pt]
\small
\textbf{Input:}~\iconev~Events \quad \iconcam~Frames \quad \icontext~Text
\quad\textbf{Repr.:}~VG = Voxel Grid,\; EF = Event Frame,\; TS = Time Surface,\; Graph = Event Graph,\; Spike = Spike Tensor,\; Learned = Learned Embedding,\; Raw = Raw Events
\quad\textbf{Arch.:}~CNN,\; RNN,\; ViT = Vision Transformer,\; GNN = Graph Neural Network,\; SNN = Spiking Neural Network,\; SSM = State Space Model,\; MoE = Mixture of Experts,\; ``--''\,=\,classical/geometric
\quad\textbf{Cat.} numbers match circled labels in \S\ref{sec:perception}: \circnum{1}\,CNN/RNN/Graph Det. \circnum{2}\,Transformer Det. \circnum{3}\,SNN/Efficient Det. \circnum{4}\,Supervised Seg. \circnum{5}\,Annotation-Efficient Seg. \circnum{6}\,Motion/Instance Seg. \circnum{7}\,Monocular Depth \circnum{8}\,Stereo Depth \circnum{9}\,Optical Flow \circnum{10}\,Frame-Event Track. \circnum{11}\,Event-Only Track. \circnum{12}\,Long-Term Track. \circnum{13}\,Action/Gesture \circnum{14}\,Body Pose \circnum{15}\,Hand/Ego Pose.
\end{tcolorbox}
\end{table*}


% ============================================================
% TABLE 2: Reconstruction Catalog
% ============================================================
\begin{table*}[p]
\centering
\caption{Dedicated event-stream and 2D reconstruction paradigms. Event-to-event methods restore the spatial resolution of the asynchronous signal itself; visual methods are then divided by sensing input into event-only appearance reconstruction and event-guided RGB restoration. Diffusion denotes a task-trained diffusion model unless the method appears in the Foundation-Prior table.}
\label{tab:recon_catalog}
\resizebox{\textwidth}{!}{
\begin{tabular}{l l c c c l}
\toprule
\textbf{Method} & \textbf{Venue} & \textbf{In} & \textbf{Out} & \textbf{Model} & \textbf{Objective} \\
\midrule
\multicolumn{6}{l}{\cellcolor{gray!8}\textit{Event-Stream Restoration (Event-to-Event)}} \\
\midrule
BMCNet~\cite{huang2024bilateral} & CVPR'24 & \iconev & HR events & CNN & Event-stream SR \\
\rowalt
NeurSR~\cite{zhang2024neuromorphic} & TCSVT'24 & \iconev & HR events & CNN & Self-supervised event SR \\
UPNSNN~\cite{xu2026ultralight} & AAAI'26 & \iconev & HR events & SNN & Efficient event-stream SR \\
\midrule
\multicolumn{6}{l}{\cellcolor{gray!8}\textit{Event-Only Image and Video Reconstruction}} \\
\midrule
E2VID~\cite{rebecq2019e2vid} & CVPR'19 & \iconev & Video & RNN & E2V \\
\rowalt
cGAN-E2V~\cite{wang2019cgane2v} & CVPR'19 & \iconev & Video & GAN & E2V \\
E2VID+~\cite{stoffregen2020hqf} & ECCV'20 & \iconev & Video & RNN & E2V \\
\rowalt
Dark-E2V~\cite{song2020darke2v} & ECCV'20 & \iconev & Video & CNN & Self-supervised E2V \\
ET-Net~\cite{weng2021etnet} & ICCV'21 & \iconev & Video & ViT & E2V \\
\rowalt
PhotoConst.~\cite{paredes2021backtoeventbasics} & CVPR'21 & \iconev & Video & CNN & Self-supervised E2V \\
EVSNN~\cite{zhu2022evsnn} & CVPR'22 & \iconev & Video & SNN & E2V \\
\rowalt
E-SAI~\cite{yu2023esai} & TPAMI'23 & \iconev & Video & CNN & E2V / see-through imaging \\
HyperE2VID~\cite{ercan2024hypere2vid} & TIP'24 & \iconev & Video & RNN & E2V \\
\rowalt
TRG-Diffusion~\cite{zhu2024trgdiffusion} & ECCV'24 & \iconev & Video & Diff. & E2V \\
EventMamba~\cite{zhang2025eventmamba} & AAAI'25 & \iconev & Video & SSM & E2V \\
\rowalt
Noise2Image~\cite{cao2025noise2image} & Optica'25 & \iconev & Image & CNN & Static-scene recovery \\
Coded-E2LF~\cite{tsuchida2026coded} & CVPR'26 & \iconev & Light field & Model-based & 4D light-field reconstruction \\
\midrule
\multicolumn{6}{l}{\cellcolor{gray!8}\textit{Event-Guided RGB Reconstruction and Restoration}} \\
\midrule
mEDI~\cite{pan2019medi} & CVPR'19 & \iconev\iconcam & Sharp video & Hybrid & Deblurring / high-rate video \\
\rowalt
eSL-Net~\cite{wang2020eslnet} & ECCV'20 & \iconev\iconcam & HR image & CNN & Super-resolution \\
TimeLens~\cite{tulyakov2021timelens} & CVPR'21 & \iconev\iconcam & Video & CNN & Frame interpolation \\
\rowalt
EFNet~\cite{sun2022efnet} & ECCV'22 & \iconev\iconcam & Sharp image & CNN & Deblurring \\
EVDI~\cite{zhang2022evdi} & CVPR'22 & \iconev\iconcam & Video & CNN & Deblurring + interpolation \\
\rowalt
CBMNet~\cite{kim2023cbmnet} & CVPR'23 & \iconev\iconcam & Video & CNN & Frame interpolation \\
\rowalt
NeurImg-HDR~\cite{han2023neurimghdr} & TPAMI'23 & \iconev\iconcam & HDR image & CNN & HDR \\
SAN~\cite{zhang2023san} & ICCV'23 & \iconev\iconcam & Sharp image & CNN & Deblurring \\
\rowalt
STIR~\cite{lu2023stir} & CVPR'23 & \iconev\iconcam & HR video & INR & Super-resolution \\
ELEDNet~\cite{kim2024elednet} & ECCV'24 & \iconev\iconcam & Video & CNN & Low-light + deblurring \\
\rowalt
EvLight~\cite{liang2024evlight} & CVPR'24 & \iconev\iconcam & Image & CNN & Low-light \\
EvTexture~\cite{kai2024evtexture} & ICML'24 & \iconev\iconcam & HR video & CNN & Super-resolution \\
\rowalt
NER-Net~\cite{liu2024rled} & CVPR'24 & \iconev\iconcam & Video & CNN & Low-light \\
TimeLens-XL~\cite{ma2024timelensxl} & ECCV'24 & \iconev\iconcam & Video & CNN & Frame interpolation \\
\rowalt
UniINR~\cite{lu2024uniinr} & ECCV'24 & \iconev\iconcam & Video & INR & RS correction + deblur + VFI \\
ClearSight~\cite{lin2025clearsight} & ICCV'25 & \iconev\iconcam & Sharp image & CNN & Deblurring \\
\rowalt
EVDM~\cite{zhang2025evdm} & ICCV'25 & \iconev\iconcam & Video & SSM & Deblurring \\
EvEnhancer~\cite{lu2025evenhancer} & CVPR'25 & \iconev\iconcam & HR video & INR & Continuous space-time SR \\
\rowalt
EventHDR~\cite{zou2025eventhdr} & TPAMI'25 & \iconev\iconcam & HDR video & CNN & HDR / high-speed imaging \\
SEE~\cite{lu2025see} & NeurIPS'25 & \iconev\iconcam & Video & CNN & Broad-range illumination \\
TimeTracker~\cite{liu2025timetracker} & CVPR'25 & \iconev\iconcam & Video & CNN & Frame interpolation \\
\rowalt
NEC-Diff~\cite{liu2026necdiff} & CVPR'26 & \iconev\iconcam & Image & Diff. & Extreme low-light RAW \\
\rowalt
Self-EHDRI~\cite{li2024hdr} & PR'26 & \iconev\iconcam & HDR image & CNN & HDR \\
\bottomrule
\end{tabular}
}
\end{table*}

\begin{table*}[p]
\centering
\caption{Dedicated 3D reconstruction paradigms. Input modality defines the primary groups; geometry, neural fields, and Gaussian Splatting are retained as representation attributes rather than competing top-level taxonomies.}
\label{tab:recon_catalog_3d}
\resizebox{\textwidth}{!}{
\begin{tabular}{l l c c c l}
\toprule
\textbf{Method} & \textbf{Venue} & \textbf{In} & \textbf{Output} & \textbf{Representation} & \textbf{Setting} \\
\midrule
\multicolumn{6}{l}{\cellcolor{gray!8}\textit{Event-Only 3D Reconstruction}} \\
\midrule
EMVS~\cite{rebecq2018emvs} & IJCV'18 & \iconev & Depth & Geometry & Multi-view stereo \\
\rowalt Semi-Dense~\cite{zhou2018semidense} & ECCV'18 & \iconev & Depth & Geometry & Stereo \\
E-NeRF~\cite{klenk2023enerf} & RA-L'23 & \iconev & Novel views & NeRF & Static \\
\rowalt EventNeRF~\cite{rudnev2023eventnerf} & CVPR'23 & \iconev & Novel views & NeRF & Static \\
Robust-e-NeRF~\cite{low2023robustenerf} & ICCV'23 & \iconev & Novel views & NeRF & Sparse/noisy events \\
\rowalt Deblur-e-NeRF~\cite{low2024deblurenerf} & ECCV'24 & \iconev & Novel views & NeRF & Blurred events \\
EvDNeRF~\cite{bhattacharya2024evdnerf} & WACV'24 & \iconev & Novel views & NeRF & Dynamic \\
\rowalt Event-3DGS~\cite{lan2024event3dgs} & NeurIPS'24 & \iconev & Novel views & 3DGS & Static \\
Event3DGS~\cite{xiong2024event3dgs} & CoRL'24 & \iconev & Novel views & 3DGS & High-speed egomotion \\
\rowalt EventPS~\cite{yu2024eventps} & CVPR'24 & \iconev & Normals & Geometry & Photometric stereo \\
PAEv3d~\cite{wang2024paev3d} & ICRA'24 & \iconev & Novel views & NeRF & Physical priors \\
\rowalt E-3DGS~\cite{zahid2025e3dgslarge} & 3DV'25 & \iconev & Novel views & 3DGS & Large/unbounded scenes \\
E-4DGS~\cite{wang2025e4dgs} & ACM MM'25 & \iconev & Dynamic views & 4DGS & Multi-view dynamic \\
\rowalt E-NeMF~\cite{wang2025enemf} & ICCV'25 & \iconev & Space-time views & Neural field & Dynamic \\
Elite-EvGS~\cite{zhang2024eliteevgs} & ICRA'25 & \iconev & Novel views & 3DGS & E2V-task prior \\
Ev4DGS~\cite{liu2025ev4dgs} & BMVC'25 & \iconev & Dynamic views & 4DGS & Monocular dynamic \\
\rowalt EventSplat~\cite{yura2025eventsplat} & CVPR'25 & \iconev & Novel views & 3DGS & Real-time rendering \\
IncEventGS~\cite{huang2025inceventgs} & CVPR'25 & \iconev & Novel views & 3DGS & Pose-free \\
\rowalt E2EGS~\cite{li2026e2egs} & CVPR'26 & \iconev & Novel views & 3DGS & Pose-free \\
EventNeuS~\cite{sachan2026eventneus} & 3DV'26 & \iconev & Mesh & NeuS & Single camera \\
\rowalt RoEL~\cite{bae2026roel} & T-RO'26 & \iconev & 3D lines & Geometry & Robust line mapping \\
\midrule
\multicolumn{6}{l}{\cellcolor{gray!8}\textit{Event-Assisted Multimodal 3D Reconstruction}} \\
\midrule
E2NeRF~\cite{qi2023e2nerf} & ICCV'23 & \iconev\iconcam & Novel views & NeRF & Blurry RGB \\
\rowalt BeNeRF~\cite{li2024benerf} & ECCV'24 & \iconev\iconcam & Novel views & NeRF & Single blurry RGB \\
E2GS~\cite{deguchi2024e2gs} & ICIP'24 & \iconev\iconcam & Novel views & 3DGS & Event-enhanced \\
\rowalt EvDeblurNeRF~\cite{cannici2024evdeblurnerf} & CVPR'24 & \iconev\iconcam & Novel views & NeRF & Deblurring \\
EvGGS~\cite{wang2024evggs} & ICML'24 & \iconev\iconcam & Novel views & 3DGS & Generalizable \\
\rowalt DynEventNeRF~\cite{rudnev2025dyneventnerf} & CVPRW'25 & \iconev\iconcam & Dynamic views & NeRF & Multi-view dynamic \\
EF-3DGS~\cite{liao2025ef3dgs} & NeurIPS'25 & \iconev\iconcam & Novel views & 3DGS & Free trajectory \\
\rowalt EGS-SLAM~\cite{chen2025egsslam} & RA-L'25 & \iconev\iconcam & Map/views & 3DGS & RGB-D SLAM \\
LSE-NeRF~\cite{liao2025lsenerf} & 3DV'25 & \iconev\iconcam & Novel views & NeRF & Sensor-error modeling \\
\rowalt HDR-NeRF-Events~\cite{qi2026seeing} & CVPR'26 & \iconev\iconcam & HDR views & NeRF & Deblur + HDR \\
\bottomrule
\end{tabular}
}
\end{table*}

\begin{table*}[p]
\centering
\caption{Foundation-Prior reconstruction. Inclusion requires a large-scale pretrained model to contribute materially through adaptation, generation, initialization, supervision, pseudo-labeling, or structured validation. Merely using a diffusion architecture, Transformer, NeRF, or 3DGS is insufficient.}
\label{tab:recon_catalog_foundation}
\resizebox{\textwidth}{!}{
\begin{tabular}{l l c l l l}
\toprule
\textbf{Method} & \textbf{Venue} & \textbf{Dim.} & \textbf{Input / Objective} & \textbf{Imported prior} & \textbf{Role in reconstruction} \\
\midrule
LaSe-E2V~\cite{chen2024lase} & NeurIPS'24 & 2D & \iconev\icontext / E2V & SD2.1 + video LDM + CLIP & Backbone initialization and semantic generation \\
\rowalt REVDM~\cite{chen2025revdm} & CVPR'25 & 2D & \iconev\iconcam / VFI & Pretrained video diffusion & Backbone repurposing \\
Semantic-E2VID~\cite{wu2025semantice2vid} & arXiv'25 & 2D & \iconev / E2V & SAM & Cross-modal distillation and semantic supervision \\
\rowalt EGVD~\cite{wang2025egvd} & ICCVW'25 & 2D & \iconev\iconcam / VFI & Stable Video Diffusion & Selective fine-tuning \\
EPA~\cite{li2025epa} & NeurIPS'25 & 2D & \iconev\iconcam / VFI & DINO & Perceptual supervisory space \\
\rowalt EVDiffuser~\cite{chen2025evdiffuser} & NeurIPS'25 & 2D & \iconev\iconcam / low-light & CogVideoX-I2V & LoRA adaptation and generation \\
EvDiff~\cite{li2025evdiff} & arXiv'25 & 2D & \iconev / E2V & Stable Diffusion 3 & One-step adaptation and surrogate training \\
\rowalt HDRev-Diff~\cite{yang2025hdrevdiff} & ICCV'25 & 2D & \iconev\iconcam / HDR & Pretrained latent diffusion & Frozen generative prior and event control \\
EvLight++~\cite{chen2026evlightpp} & TPAMI'26 & 2D & \iconev\iconcam / low-light & SAM + depth foundation model & Pseudo-labeling and structured validation \\
\rowalt IE2Video~\cite{torbunov2025ie2video} & arXiv'25 & 2D & \iconev\iconcam / E2V & LTX-Video & Event adapter and LoRA \\
UniE2F~\cite{xu2026unie2f} & arXiv'26 & 2D & \iconev / E2V, VFI, prediction & Pretrained video diffusion & Generative prior and sampling control \\
\rowalt DiET-GS~\cite{lee2025dietgs} & CVPR'25 & 3D & \iconev\iconcam / deblurred NVS & Pretrained image diffusion & Reconstruction regularization \\
DeblurSplat~\cite{zhang2025deblursplat} & TMM'26 & 3D & \iconev\iconcam / deblurred NVS & DUSt3R & Geometric initialization \\
\bottomrule
\end{tabular}
}
\vspace{2pt}
\begin{tcolorbox}[colback=gray!5, colframe=gray!50, boxrule=0.4pt, arc=2pt, left=4pt, right=4pt, top=2pt, bottom=2pt]
\small
\textbf{Input:}~\iconev~events, \iconcam~RGB/frame input, \icontext~text.
Task-trained diffusion without a large-scale pretrained model (TRG-Diffusion and NEC-Diff) remains in \cref{tab:recon_catalog}; E2V teachers (Elite-EvGS and EventSplat) are task-pretrained priors rather than foundation models. Event-to-event super-resolution is retained as a distinct signal-restoration block in \cref{tab:recon_catalog}.
\end{tcolorbox}
\end{table*}

% ============================================================
% TABLE 3: Understanding Catalog
% ============================================================
\begin{table*}[p]
\centering
\caption{Comprehensive catalog of event camera understanding methods leveraging large-scale pre-trained models, self-supervised learning, simulation, and generation. Methods are grouped by research theme and listed chronologically within each group.}
\label{tab:understanding_catalog}
\resizebox{\textwidth}{!}{
\begin{tabular}{rl l c c l l c}
\toprule
\textbf{\#} & \textbf{Method} & \textbf{Venue} & \textbf{Foundation} & \textbf{In} & \textbf{Technique} & \textbf{Task} & \textbf{Theme} \\
\midrule
\multicolumn{8}{l}{\cellcolor{gray!8}\textit{4.1 Open-Vocabulary Perception}} \\
\midrule
%% --- Image-Level Recognition ---
1 & EventCLIP~\cite{wu2023eventclip} & arXiv'23 & CLIP & \iconev\icontext & Prompt tuning for events & Recognition & \circnum{1} \\
\rowalt
2 & ExACT~\cite{zhou2024exact} & CVPR'24 & CLIP & \iconev\icontext & Language-guided reasoning & Action Recog. & \circnum{1} \\
3 & EventBind~\cite{kim2024eventbind} & ECCV'24 & CLIP & \iconev\iconcam\icontext & Unified multi-modal binding & Open-world & \circnum{1} \\
\rowalt
4 & CEIA~\cite{zhou2024ceia} & ECCVW'24 & CLIP & \iconev\iconcam\icontext & Contrastive event-image align. & Open-world & \circnum{1} \\
5 & Expanding Event~\cite{zhang2024expandingevent} & arXiv'24 & CLIP & \iconev\icontext & Robust CLIP encoder & Recognition & \circnum{1} \\
\rowalt
6 & EZSR~\cite{cho2025ezsr} & CVPR'25 & CLIP & \iconev\icontext & Zero-shot recognition & Recognition & \circnum{1} \\
%% --- Dense Segmentation & Detection ---
7 & OpenESS~\cite{kong2024openess} & CVPR'24 & CLIP & \iconev\iconcam\icontext & Open-vocab. semantic seg. & Seg. & \circnum{2} \\
\rowalt
8 & EventSAM~\cite{chen2024eventsam} & CVPR'24 & SAM & \iconev & Pivotal token adaptation & Seg. & \circnum{2} \\
9 & SAM-E-Adapter~\cite{li2024sameventadapter} & ICRA'24 & SAM & \iconev\iconcam & Adapter for SAM & Seg. & \circnum{2} \\
\rowalt
10 & OVOSE~\cite{kim2024ovose} & arXiv'24 & CLIP & \iconev\icontext & Open-vocab. event seg. & Seg. & \circnum{2} \\
11 & OV-Detection~\cite{zhang2025ovdetection} & arXiv'25 & CLIP & \iconev\icontext & Adaptive event slicing & Det. & \circnum{2} \\
\rowalt
12 & EvSAM~\cite{li2025evsam} & TOMM'25 & SAM & \iconev\iconcam & Event-assisted SAM & Seg. & \circnum{2} \\
13 & SEAL~\cite{li2026seal} & ICLR'26 & CLIP & \iconev\icontext & Language-guided seg. & Seg. & \circnum{2} \\
%% --- Self-Supervised Pre-training ---
\rowalt
14 & ECDP~\cite{yang2023ecdp} & ICCV'23 & None & \iconev & Masked image modeling & Pre-training & \circnum{3} \\
15 & MEM~\cite{klenk2024mem} & WACV'24 & None & \iconev & Masked event modeling & Pre-training & \circnum{3} \\
\rowalt
16 & DECP~\cite{li2024decp} & arXiv'24 & None & \iconev & Disentangled masked model. & Pre-training & \circnum{3} \\
17 & ECDP-Dense~\cite{wu2024ecdpdense} & ECCV'24 & None & \iconev & Dense pre-training & Pre-training & \circnum{3} \\
\rowalt
18 & SD2Event~\cite{gu2024sd2event} & CVPR'24 & None & \iconev & Dynamic detector SSL & Pre-training & \circnum{3} \\
19 & EvRepSL~\cite{zhang2024evrepsl} & TIP'24 & None & \iconev & Self-supervised repr. & Pre-training & \circnum{3} \\
\rowalt
20 & TESPEC~\cite{wang2025tespec} & ICCV'25 & None & \iconev & Temporal-enhanced SSL & Pre-training & \circnum{3} \\
21 & STP~\cite{liu2025stp} & ICCV'25 & None & \iconev & Adaptive prompt fusion & Pre-training & \circnum{3} \\
\rowalt
22 & EventPretrain~\cite{zhang2025eventpretrain} & ACM MM'25 & None & \iconev & Physics-inspired SSL & Pre-training & \circnum{3} \\
23 & CM3AE~\cite{wang2025cm3ae} & ACM MM'25 & None & \iconev\iconcam & Unified RGB-event pretraining & Pre-training & \circnum{3} \\
\rowalt
24 & GEP & CVPR'26 & None & \iconev & Generative event pretrain. & Pre-training & \circnum{3} \\
%% --- Cross-Modal Transfer ---
25 & EventDance~\cite{zheng2024eventdance} & CVPR'24 & None & \iconev & Source-free cross-modal adapt. & Recognition & \circnum{4} \\
\rowalt
26 & S5-ViT~\cite{zubic2024s5vit} & CVPR'24 & None & \iconev & State space model transfer & Multiple & \circnum{4} \\
27 & Spike-DINOv2~\cite{zhang2024spikedinov2} & arXiv'24 & DINOv2 & \iconev & Cross-modal distillation & Depth & \circnum{4} \\
\rowalt
28 & VELoRA & arXiv'24 & None & \iconev\iconcam & Low-rank RGB-event adapt. & Recognition & \circnum{4} \\
29 & EventFly~\cite{kong2025eventfly} & CVPR'25 & None & \iconev\iconcam & Frame-to-event distill. & Multiple & \circnum{4} \\
\rowalt
30 & Depth AnyEvent~\cite{wu2025eventdam} & ICCV'25 & Depth Any. & \iconev & Dense-to-sparse distill. & Depth & \circnum{4} \\
31 & Semantic-E2VID~\cite{wu2025semantice2vid} & arXiv'25 & SAM & \iconev & Semantic-enriched E2V & Recon. & \circnum{4} \\
\rowalt
32 & TGVFM~\cite{chen2025tgvfm} & arXiv'25 & Depth Any. & \iconev & Temporal-guided VFM & Multiple & \circnum{4} \\
33 & ADAE & arXiv'25 & Depth Any. & \iconev & Depth adaptation w/ events & Depth & \circnum{4} \\
\rowalt
34 & FFEvent~\cite{li2026ffevent} & ESA'26 & None & \iconev & Fourier-based transfer & Recognition & \circnum{4} \\
\midrule
\multicolumn{8}{l}{\cellcolor{gray!8}\textit{4.2 Scene Understanding \& Reasoning}} \\
\midrule
%% --- Event-Based MLLMs ---
35 & LLM Zero-Shot~\cite{yang2025llmzeroshot} & ICASSP'25 & GPT-4o & \iconev\icontext & Zero-shot text prompt & Recognition & \circnum{5} \\
\rowalt
36 & EventGPT~\cite{wang2025eventgpt} & CVPR'25 & LLaVA & \iconev\iconcam\icontext & 3-stage MLLM training & Understanding & \circnum{5} \\
37 & LLM-EvRep~\cite{zhang2025llmevrep} & WWW'25 & LLM & \iconev\icontext & Self-supervised repr. for LLM & Recognition & \circnum{5} \\
\rowalt
38 & EventVL~\cite{chen2025eventvl} & arXiv'25 & LLaVA & \iconev\icontext & Event-text alignment MLLM & Understanding & \circnum{5} \\
39 & EventBench & arXiv'25 & MLLM & \iconev\icontext & MLLM benchmark & Understanding & \circnum{5} \\
\rowalt
40 & EvQA & arXiv'25 & MLLM & \iconev\icontext & Recon.-bridge VQA & Understanding & \circnum{5} \\
41 & EventFlash~\cite{li2026eventflash} & arXiv'26 & MLLM & \iconev\icontext & Efficient event MLLM & Understanding & \circnum{5} \\
\rowalt
42 & Event-MLLM & arXiv'26 & MLLM & \iconev\icontext & Illumination-robust MLLM & Understanding & \circnum{5} \\
%% --- Visual Grounding & Embodied Intelligence ---
43 & NeuroViG~\cite{wang2025neurovig} & WACV'25 & None & \iconev\iconcam\icontext & Event-enhanced grounding & Grounding & \circnum{6} \\
\rowalt
44 & EP-VLM~\cite{chen2025epvlm} & arXiv'25 & VLM & \iconev\icontext & Event-priori VLM & Understanding & \circnum{6} \\
45 & Talk2Event~\cite{kong2025talk2event} & NeurIPS'25 & -- & \iconev\icontext & Grounded scene understanding & Grounding & \circnum{6} \\
\rowalt
46 & LLaFEA & arXiv'25 & LMM & \iconev\iconcam\icontext & Frame-event complementary & Understanding & \circnum{6} \\
47 & E-VLA~\cite{wang2026evla} & arXiv'26 & VLA & \iconev\iconcam\icontext & Event-augmented VLA & Robotics & \circnum{6} \\
%% --- Language-Guided Recognition ---
\rowalt
48 & Ev-LaFOR~\cite{cho2023evlafor} & ICCV'23 & None & \iconev & Joint recon. \& recog. & Recognition & \circnum{7} \\
49 & EventDance++~\cite{zheng2024eventdancepp} & arXiv'24 & None & \iconev\icontext & Language-guided adaptation & Recognition & \circnum{7} \\
\rowalt
50 & ExACT~\cite{zhou2024exact} & CVPR'24 & CLIP & \iconev\icontext & Language-guided reasoning & Action Recog. & \circnum{7} \\
51 & ESTR-CoT & arXiv'25 & LLM & \iconev\icontext & Chain-of-thought reasoning & Text Recog. & \circnum{7} \\
\midrule
\multicolumn{8}{l}{\cellcolor{gray!8}\textit{4.3 Event Data Simulation \& Generation}} \\
\midrule
%% --- Physics-Based Simulators ---
52 & ESIM~\cite{rebecq2018esim} & CoRL'18 & -- & \iconcam & Physics-based event sim. & Simulation & \circnum{8} \\
\rowalt
53 & rpg\_vid2e~\cite{gehrig2020vid2e} & CVPR'20 & -- & \iconcam & Video-to-event conversion & Simulation & \circnum{8} \\
54 & v2e~\cite{hu2021v2e} & CVPRW'21 & -- & \iconcam & Realistic DVS simulation & Simulation & \circnum{8} \\
\rowalt
55 & DVS-Voltmeter~\cite{lin2022dvsvoltmeter} & ECCV'22 & -- & \iconcam & Stochastic event sim. & Simulation & \circnum{8} \\
56 & ADV2E & arXiv'24 & -- & \iconcam & Analogue circuit sim. & Simulation & \circnum{8} \\
\rowalt
57 & SENPI & arXiv'25 & -- & \iconcam & PyTorch event gen. tool & Simulation & \circnum{8} \\
%% --- Learned & Neural Generation ---
58 & Text-to-Events~\cite{yang2024text2events} & NICE'24 & SD & \icontext & Text-conditioned synthesis & Generation & \circnum{9} \\
\rowalt
59 & ShapeAug & arXiv'24 & None & \iconev & Shape-based augmentation & Generation & \circnum{9} \\
60 & ControlEvents~\cite{li2025controlevents} & WACV'26 & SD & \icontext\iconcam & Controllable event gen. & Generation & \circnum{9} \\
%% --- Sim-to-Real Transfer ---
\rowalt
61 & E2VID+~\cite{stoffregen2020hqf} & ECCV'20 & -- & \iconev & Reduce sim-to-real gap & Sim2Real & \circnum{10} \\
62 & DA4Event & arXiv'21 & -- & \iconev & Domain adapt. for events & Sim2Real & \circnum{10} \\
\rowalt
63 & N-ROD & arXiv'21 & -- & \iconev & Syn-to-real domain adapt. & Sim2Real & \circnum{10} \\
65 & EQS & CVPRW'25 & -- & \iconev & Event realism scoring & Sim2Real & \circnum{10} \\
66 & CARLA-DVS & arXiv'25 & -- & \iconev & Sim-to-real gap analysis & Sim2Real & \circnum{10} \\
\bottomrule
\end{tabular}
}

\vspace{2pt}
\begin{tcolorbox}[colback=gray!5, colframe=gray!50, boxrule=0.4pt, arc=2pt, left=4pt, right=4pt, top=2pt, bottom=2pt]
\small
\textbf{Input:}~\iconev~Events \quad \iconcam~Frames \quad \icontext~Text
\quad\textbf{Foundation:}~CLIP,\; SAM,\; DINOv2,\; Depth Any.\,=\,Depth Anything,\; LLaVA,\; SD\,=\,Stable Diffusion,\; VLM\,=\,Vision-Language Model,\; VLA\,=\,Vision-Language-Action,\; MLLM\,=\,Multimodal LLM,\; LMM\,=\,Large Multimodal Model,\; None\,=\,no foundation model,\; ``--''\,=\,not applicable
\quad\textbf{Theme:} \circnum{1}\,Image-Level Recog. \circnum{2}\,Dense Seg.\,\&\,Det. \circnum{3}\,Self-Supervised Pre-training \circnum{4}\,Cross-Modal Transfer \circnum{5}\,MLLMs \circnum{6}\,Grounding\,\&\,Embodied \circnum{7}\,Language-Guided \circnum{8}\,Physics-Based Sim. \circnum{9}\,Learned/Neural Gen. \circnum{10}\,Sim-to-Real.
\end{tcolorbox}
\end{table*}



% ======================================================================
% 3.2 Reconstruction
% ======================================================================
\subsection{Event Camera Reconstruction}
\label{sec:reconstruction}

Event camera reconstruction restores or derives dense signals---enhanced event streams, images, videos, light fields, geometry, and view-synthesis representations---from asynchronous brightness changes, optionally aided by conventional frames. A taxonomy based only on network families is unstable: a Transformer or diffusion architecture does not reveal what is reconstructed, what measurements are available, or whether knowledge was imported from large-scale pretraining. We therefore use two top-level paradigms: \textit{Dedicated Reconstruction Paradigms}, whose reconstruction knowledge is learned or optimized for the event task, and \textit{Foundation-Prior Reconstruction}, in which a large-scale pretrained model materially contributes through reconstruction, training, data enrichment, or structured validation. Within the dedicated branch, the output domain first distinguishes event-stream restoration from visual reconstruction; visual methods are then organized by input modality, task objective, and 3D representation. \Cref{tab:recon_catalog,tab:recon_catalog_3d,tab:recon_catalog_foundation} catalog 91 in-scope methods under this scheme.

\subsubsection{Dedicated Reconstruction Paradigms}
\label{sec:recon_dedicated}

We use ``dedicated'' as a neutral operational term: these methods may generalize across sensors or scenes, but their core prior is optimized specifically for event reconstruction. This avoids the misleading implication that every diffusion, Transformer, NeRF, or 3DGS method belongs to the large-model branch.

\paragraph{Event-Stream Restoration.}
This small but distinct line reconstructs a higher-resolution asynchronous signal rather than an intensity image. BMCNet~\cite{huang2024bilateral} separates positive and negative streams and exchanges complementary spatio-temporal information; NeurSR~\cite{zhang2024neuromorphic} performs self-supervised, per-input neuromorphic super-resolution without prior training; and UPNSNN~\cite{xu2026ultralight} uses a polarity-split lightweight SNN to preserve asynchronous output for efficient deployment. These methods do not fit perception, because they infer no semantic or geometric target, and they should not be merged with event-to-image super-resolution. We therefore retain them as an event-to-event restoration block within reconstruction.

\paragraph{2D Event-Only Image and Video Reconstruction.}
Event-only reconstruction must infer absolute appearance from relative brightness changes, making static texture and color intrinsically ambiguous. E2VID~\cite{rebecq2019e2vid} established the recurrent encoder--decoder baseline, while E2VID+~\cite{stoffregen2020hqf} reduced the simulation-to-real gap. Photometric-constancy training~\cite{paredes2021backtoeventbasics} showed that an event formation model can replace paired ground truth. Later work changed the temporal operator rather than the task definition: ET-Net~\cite{weng2021etnet} adopted attention, EVSNN~\cite{zhu2022evsnn} used spiking dynamics, HyperE2VID~\cite{ercan2024hypere2vid} generated adaptive weights, and EventMamba~\cite{zhang2025eventmamba} used state-space scanning. Task-trained generative models, including cGAN-E2V~\cite{wang2019cgane2v} and TRG-Diffusion~\cite{zhu2024trgdiffusion}, improve perceptual realism but do not import a foundation model merely by using GAN or diffusion objectives. Event-only sensing also supports specialized computational imaging: Noise2Image~\cite{cao2025noise2image} exploits sensor noise for static-scene recovery, whereas Coded-E2LF~\cite{tsuchida2026coded} reconstructs a 4D light field using coded optics.

\paragraph{2D Event-Guided RGB Reconstruction and Restoration.}
Here an RGB/LDR/blurred frame supplies appearance and color, while events supply exposure-time motion and high-dynamic-range change. The task objective gives the clearest subdivision. For \textit{frame interpolation}, TimeLens~\cite{tulyakov2021timelens}, CBMNet~\cite{kim2023cbmnet}, TimeLens-XL~\cite{ma2024timelensxl}, and TimeTracker~\cite{liu2025timetracker} progressively address synthesis, bidirectional motion, large displacement, and non-linear trajectories. For \textit{motion deblurring}, mEDI~\cite{pan2019medi} introduced event-based exposure inversion, EFNet~\cite{sun2022efnet} learned cross-modal fusion, and EVDM~\cite{zhang2025evdm} targets real-world video with state-space modeling. Unified methods are marked explicitly rather than forced into one task: EVDI~\cite{zhang2022evdi} combines deblurring and interpolation, and UniINR~\cite{lu2024uniinr} additionally handles rolling-shutter correction through a coordinate-based representation.

Spatial and radiometric restoration form three further objectives. Event-guided super-resolution progressed from eSL-Net~\cite{wang2020eslnet} to spatial-temporal implicit representations~\cite{lu2023stir}, texture-aware fusion~\cite{kai2024evtexture}, and continuous space-time recovery~\cite{lu2025evenhancer}. HDR methods fuse the event sensor's dynamic range with frame color, from NeurImg-HDR~\cite{han2023neurimghdr} and EventHDR~\cite{zou2025eventhdr} to broad-range illumination adjustment~\cite{lu2025see}. Low-light restoration is represented by NER-Net~\cite{liu2024rled}, EvLight~\cite{liang2024evlight}, and the event--RAW NEC-Diff model~\cite{liu2026necdiff}.

\paragraph{3D Event-Only Reconstruction.}
\label{sec:recon3d}
The event-only line progresses in representation while retaining the harder sensing condition. Geometric methods such as EMVS~\cite{rebecq2018emvs} and stereo semi-dense reconstruction~\cite{zhou2018semidense} accumulate event consistency into depth; EventPS~\cite{yu2024eventps} instead recovers surface normals through photometric stereo. EventNeRF~\cite{rudnev2023eventnerf} and E-NeRF~\cite{klenk2023enerf} introduced neural radiance fields learned directly from moving event cameras, followed by robustness to sparse/noisy events~\cite{low2023robustenerf}, dynamic radiance fields~\cite{bhattacharya2024evdnerf}, neural surfaces~\cite{sachan2026eventneus}, and neural motion fields~\cite{wang2025enemf}.

Since 2024, 3DGS has shifted the emphasis toward fast rendering and joint pose estimation. Event3DGS~\cite{xiong2024event3dgs} targets high-speed robot egomotion, while Event-3DGS~\cite{lan2024event3dgs} establishes event-only Gaussian reconstruction as a general formulation. IncEventGS~\cite{huang2025inceventgs} and E2EGS~\cite{li2026e2egs} remove externally supplied poses; E-3DGS~\cite{zahid2025e3dgslarge} scales to large and unbounded scenes; Ev4DGS~\cite{liu2025ev4dgs} and E-4DGS~\cite{wang2025e4dgs} extend reconstruction to non-rigid or multi-view dynamics. Elite-EvGS~\cite{zhang2024eliteevgs} and EventSplat~\cite{yura2025eventsplat} use pretrained E2V networks, but those are task-specific reconstruction teachers rather than general foundation models.

\paragraph{3D Event-Assisted Multimodal Reconstruction.}
RGB or RGB-D inputs alleviate the missing-appearance problem and make events a physical cue for blur, exposure, trajectory, or temporal change. In NeRF, E2NeRF~\cite{qi2023e2nerf}, EvDeblurNeRF~\cite{cannici2024evdeblurnerf}, and BeNeRF~\cite{li2024benerf} reconstruct from degraded frames, while HDR-NeRF-Events~\cite{qi2026seeing} targets high-dynamic-range radiance fields. In 3DGS, E2GS~\cite{deguchi2024e2gs} and EF-3DGS~\cite{liao2025ef3dgs} exploit event cues for quality or free-trajectory capture; EvGGS~\cite{wang2024evggs} studies feed-forward generalization, and EGS-SLAM~\cite{chen2025egsslam} incorporates events into RGB-D Gaussian SLAM. These representation labels remain essential table attributes, but input modality better exposes the information available to each method.

\subsubsection{Foundation-Prior Reconstruction}
\label{sec:recon_foundation}

Foundation-Prior Reconstruction requires an external, large-scale pretrained model to contribute materially to the method. Its role may be backbone adaptation, generative regularization, geometric initialization, distillation, supervision, pseudo-labeling, or structured evaluation that exposes capabilities of the reconstructed output. This criterion still excludes a diffusion model trained from scratch for one event task, ordinary ImageNet initialization, and an E2V warm start.

\paragraph{2D Foundation-Prior Reconstruction.}
LaSe-E2V~\cite{chen2024lase} initializes a language-conditioned video latent diffusion model from Stable Diffusion and large-scale video pretraining, then injects event structure and CLIP text semantics. REVDM~\cite{chen2025revdm} directly repurposes a pretrained video diffusion model for event-guided interpolation; EGVD~\cite{wang2025egvd} selectively adapts Stable Video Diffusion for large motion; IE2Video~\cite{torbunov2025ie2video} injects events and a sparse RGB keyframe into LTX-Video through learned adapters and LoRA; and UniE2F~\cite{xu2026unie2f} uses a pretrained video diffusion prior for E2V, interpolation, and prediction.

Foundation priors need not be generative backbones. Semantic-E2VID~\cite{wu2025semantice2vid} distills SAM features and uses SAM-derived semantic supervision, while EPA~\cite{li2025epa} places DINO features at the center of perceptually aligned VFI supervision. EVDiffuser~\cite{chen2025evdiffuser} fine-tunes CogVideoX-I2V with modality-agnostic event/RGB conditioning. EvDiff~\cite{li2025evdiff} explicitly adopts Stable Diffusion 3 as its base model, converts it to one-step event reconstruction, and bridges scarce paired event data to Places365 through surrogate training and EvEncoder distillation.

The large-model contribution can also enter through a broader experimental pipeline. HDRev-Diff~\cite{yang2025hdrevdiff} keeps a pretrained latent diffusion model fixed and learns event-conditioned control for HDR generation. EvLight++~\cite{chen2026evlightpp} uses SAM and a monocular depth foundation model to enrich its large real-world dataset with semantic and depth pseudo-labels and to evaluate whether enhanced videos recover useful scene structure. Although these roles differ from adapting a generative backbone, both methods derive a substantive part of their contribution from large-scale pretrained knowledge and therefore belong to the large-model-era branch.

\paragraph{3D Foundation-Prior Reconstruction.}
Current evidence supports a smaller but conceptually clear branch. DiET-GS~\cite{lee2025dietgs} uses the natural-image distribution of a pretrained diffusion model to regularize event-assisted deblurring 3DGS. DeblurSplat~\cite{zhang2025deblursplat} imports DUSt3R's learned multiview geometry to initialize point clouds before event-supervised Gaussian optimization. They differ from event 3DGS methods that merely adopt the 3DGS representation: the imported pretrained knowledge participates directly in solving the scene reconstruction.

\textit{Discussion.} The new taxonomy separates two independent developments that the former architecture-based organization conflated. Dedicated models continue to advance through better physical objectives, temporal operators, and scene representations; Foundation-Prior methods draw on appearance, semantics, video dynamics, or geometry learned at scale, whether inside the reconstructor or elsewhere in the method's substantive pipeline. The latter currently remains much smaller, especially in 3D, so subdividing it by every restoration task would create sparse and unstable categories. The tables therefore retain task, input, representation, pose, and scene setting as attributes while the prose records how each pretrained model contributes. Quantitative comparisons remain in \cref{sec:appendix}.


% ======================================================================
% 3.3 Understanding
% ======================================================================
\subsection{Event Camera Understanding}
\label{sec:understanding}

The most distinctive contribution of this survey is the systematic review of event camera understanding --- an emerging area where large-scale pre-trained models are adapted to event data for semantic reasoning, open-vocabulary recognition, and language-grounded perception. Unlike the perception tasks in \cref{sec:perception}, which rely on task-specific architectures trained on event-specific labeled data, understanding methods leverage the rich visual and semantic priors encoded in foundation models (\eg, CLIP, SAM, GPT) to overcome the fundamental bottleneck of limited labeled event data. We organize this section around three themes that collectively define the frontier of event-based vision in the era of large models. \cref{fig:understanding_overview} maps the landscape of foundation model adaptation for event cameras.

\begin{figure*}[t]
\centering
\begin{tikzpicture}[
    font=\sffamily\scriptsize,
    fmbox/.style={rounded corners=3pt, draw=#1!80, fill=#1!12, thick, minimum width=2cm, minimum height=0.7cm, align=center, font=\sffamily\scriptsize\bfseries},
    methodbox/.style={rounded corners=2pt, draw=#1!60, fill=#1!6, minimum width=1.8cm, minimum height=0.5cm, align=center, font=\sffamily\tiny, inner sep=2pt},
    taskbox/.style={rounded corners=2pt, draw=gray!50, fill=gray!5, minimum width=1.8cm, minimum height=0.5cm, align=center, font=\sffamily\tiny\itshape, inner sep=2pt},
    arr/.style={-Stealth, thick, #1!60},
    larr/.style={-Stealth, thin, #1!40},
]

% Foundation Models (top row)
\node[fmbox=blue] (clip) at (0, 0) {CLIP};
\node[fmbox=green!60!black] (sam) at (3.5, 0) {SAM};
\node[fmbox=red!70!black] (diff) at (7, 0) {Diffusion\\Models};
\node[fmbox=violet] (depth) at (10.5, 0) {Depth\\Anything};
\node[fmbox=orange] (llm) at (14, 0) {LLMs /\\MLLMs};

% Title
\node[font=\sffamily\small\bfseries, text=gray!70] at (7, 1.2) {Pre-trained Foundation Models};
\draw[gray!30, thick, dashed] (-1.5, 0.7) -- (15.5, 0.7);

% Event-specific methods (middle row)
\node[methodbox=blue] (eclip) at (-0.9, -1.5) {EventCLIP};
\node[methodbox=blue] (ebind) at (0.9, -1.5) {EventBind};
\node[methodbox=blue] (ezsr) at (0, -2.3) {EZSR};

\node[methodbox=green!60!black] (esam) at (2.8, -1.5) {EventSAM};
\node[methodbox=green!60!black] (evsam) at (4.2, -1.5) {EvSAM};

\node[methodbox=red!70!black] (revdm) at (6.2, -1.5) {REVDM};

\node[methodbox=violet] (edam) at (10.5, -1.5) {EventDAM};

\node[methodbox=orange] (egpt) at (13.1, -1.5) {EventGPT};
\node[methodbox=orange] (evl) at (14.9, -1.5) {EventVL};
\node[methodbox=orange] (efl) at (14, -2.3) {EventFlash};

% Arrows from foundation models to methods
\draw[arr=blue] (clip) -- (eclip);
\draw[arr=blue] (clip) -- (ebind);
\draw[arr=green!60!black] (sam) -- (esam);
\draw[arr=green!60!black] (sam) -- (evsam);
\draw[arr=red!70!black] (diff) -- (revdm);
\draw[arr=violet] (depth) -- (edam);
\draw[arr=orange] (llm) -- (egpt);
\draw[arr=orange] (llm) -- (evl);

% Tasks (bottom row)
\node[taskbox] (t1) at (-0.5, -3.5) {Zero-shot\\Recognition};
\node[taskbox] (t2) at (2, -3.5) {Open-Vocab\\Segmentation};
\node[taskbox] (t3) at (4.5, -3.5) {Segment\\Anything};
\node[taskbox] (t4) at (7, -3.5) {Generative\\Reconstruction};
\node[taskbox] (t5) at (9.5, -3.5) {Depth\\Estimation};
\node[taskbox] (t6) at (12, -3.5) {VQA \&\\Captioning};
\node[taskbox] (t7) at (14.5, -3.5) {Embodied\\Intelligence};

% Arrows from methods to tasks
\draw[larr=blue] (ezsr) -- (t1);
\draw[larr=blue] (eclip) -- (t2);
\draw[larr=green!60!black] (esam) -- (t3);
\draw[larr=red!70!black] (revdm) -- (t4);
\draw[larr=violet] (edam) -- (t5);
\draw[larr=orange] (egpt) -- (t6);
\draw[larr=orange] (efl) -- (t7);

% Additional cross-references
\node[methodbox=blue] (openess) at (2, -2.3) {OpenESS};
\draw[larr=blue] (openess) -- (t2);

% Label rows
\node[font=\sffamily\tiny\bfseries, text=gray!50, rotate=90] at (-2, -1.5) {Methods};
\node[font=\sffamily\tiny\bfseries, text=gray!50, rotate=90] at (-2, -3.5) {Tasks};

\end{tikzpicture}
\caption{Landscape of foundation model adaptation for event camera understanding. Five families of pre-trained models (\textcolor{blue!80}{CLIP}, \textcolor{green!60!black}{SAM}, \textcolor{red!60!black}{Diffusion}, \textcolor{violet}{Depth Anything}, \textcolor{orange!80}{LLMs/MLLMs}) are adapted to event data through contrastive alignment, adapter modules, knowledge distillation, or end-to-end training. Each pathway enables qualitatively new capabilities --- from zero-shot recognition to visual question answering --- that were previously unavailable for event cameras.}
\label{fig:understanding_overview}
\end{figure*}


% ----------------------------------------------------------------------
% 3.3.1 Open-Vocabulary Perception
% ----------------------------------------------------------------------
\subsubsection{Open-Vocabulary Perception}
\label{sec:understanding_ov}

A defining capability of foundation-model-adapted event vision is \textit{open-vocabulary understanding}: recognizing or segmenting objects described by arbitrary text queries, without being limited to a predefined set of categories. This subsection consolidates work on CLIP-based alignment, open-vocabulary segmentation and detection, foundation model adaptation, self-supervised pre-training, and cross-modal knowledge transfer --- all of which share the common goal of enabling event models to generalize beyond closed-set categories.

\paragraph{CLIP-Based Alignment and Zero-Shot Recognition.}
CLIP~\cite{radford2021clip} provides a powerful image-text embedding space that can serve as a bridge between event data and natural language. EventCLIP~\cite{wu2023eventclip} was the first to adapt CLIP for event-based recognition, converting events to 2D grid representations and introducing a feature adapter for temporal aggregation. EventBind~\cite{zhou2024eventbind} (ECCV'24) proposed Hierarchical Triple Contrastive Alignment (HTCA) across image, text, and event modalities, achieving state-of-the-art on N-Caltech101 and N-ImageNet. CEIA~\cite{wei2024ceia} used images as a bridge to align events with CLIP's space via LoRA-based event encoder training, preserving CLIP's zero-shot ability while adapting to events. The Expanding Event encoder~\cite{zhou2024expanding} extended alignment to five modalities (image, event, text, sound, depth) with anti-catastrophic-forgetting mechanisms.

EZSR~\cite{cho2025ezsr} (CVPR'25) represents a breakthrough: it develops an event encoder without relying on reconstruction networks, synthesizes training events from static RGB images at scale, and achieves 47.84\% zero-shot accuracy on N-ImageNet with ViT-B/16 --- \textit{surpassing even supervised methods}. This result demonstrates that event-text alignment, when properly designed, can be more effective than event-specific supervised training.

ExACT~\cite{tang2024exact} (CVPR'24) applies CLIP-based reasoning to action recognition, using GPT-4-generated action captions to build the SeAct dataset and achieving 94.83\% accuracy on PAF through language-guided conceptual reasoning.

\paragraph{Open-Vocabulary Segmentation and Detection.}
OpenESS~\cite{kong2024openess} (CVPR'24 Highlight) was the first open-vocabulary event-based segmentation method. It transfers CLIP knowledge to event streams via two mechanisms: (i)~frame-to-event contrastive distillation, which aligns event features with CLIP image features, and (ii)~text-to-event semantic consistency regularization. Remarkably, OpenESS achieves 53.93\% mIoU on DDD17 and 43.31\% on DSEC-Semantic \textit{without any event or frame labels during training} --- demonstrating that language supervision alone can enable dense event perception.

OVOSE~\cite{bendali2024ovose} extended open-vocabulary segmentation using synthetic event data and foundation model distillation. An adaptive event stream slicing method~\cite{li2025ovdetection} brought open-vocabulary capabilities to object detection, using SNN-based slicing with CLIP distillation from an RGB teacher. SEAL~\cite{li2026seal} (ICLR'26) represents the current state of the art, enabling Semantic-aware Segment Any Events at both instance and part-level granularity through open-vocabulary event instance segmentation (OV-EIS), with four benchmarks spanning coarse-to-fine semantic granularity.

\paragraph{Foundation Model Adaptation.}
Beyond CLIP, other foundation models --- particularly the Segment Anything Model (SAM)~\cite{kirillov2023sam} and DINOv2 --- have been adapted to event data through lightweight adapters and knowledge distillation.

\noindent\circnum{1}~\textbf{SAM Adaptation.} EventSAM~\cite{chen2024eventsam} (CVPR'24) adapts SAM for event segmentation through multi-scale feature distillation that aligns event token embeddings with RGB counterparts. It also contributes the RGBE-SEG dataset. SAM-Event-Adapter~\cite{zhou2024sameventadapter} (ICRA'24) introduces a lightweight adapter that incorporates event data into SAM in a patch-compatible format. EvSAM~\cite{hu2025evsam} leverages SAM's generalization for RGB-event multimodal segmentation under adverse conditions.

\noindent\circnum{2}~\textbf{DINOv2 and Visual Foundation Models.} A spike-driven Transformer~\cite{liu2024spikedinov2} uses DINOv2 as a teacher for cross-modality knowledge distillation, achieving 49\% error reduction and 70\% energy savings for event-based depth estimation. TGVFM~\cite{zhang2025tgvfm} integrates Visual Foundation Models with temporal context fusion via Long-Range Temporal Attention and Dual Spatiotemporal Attention, achieving 16--21\% improvement over prior methods across segmentation, depth, and detection.

\noindent\circnum{3}~\textbf{Semantic Transfer for Reconstruction.} Semantic-E2VID~\cite{wu2025semantice2vid} transfers SAM's visual semantic knowledge to the event encoder for semantic-aware event-to-video reconstruction, demonstrating that foundation model priors benefit not only understanding but also reconstruction quality.

\paragraph{Self-Supervised Pre-training and Cross-Modal Transfer.}
Self-supervised pre-training learns general-purpose event representations from unlabeled data, reducing dependence on expensive annotations, while cross-modal transfer bridges the domain gap between richly annotated RGB data and sparsely annotated event data.

\noindent\circnum{1}~\textbf{Masked Modeling.} MEM~\cite{klenk2024mem} (WACV'24) trained a discrete VAE to compress event histograms into visual tokens and pre-trained a ViT with 50\% patch masking, improving classification on N-ImageNet by 7.96\%. DECP~\cite{chen2024decp} proposed data-efficient pre-training with semantic-uniform masking and disentangled reconstruction (local spatio-temporal and global semantic branches), converging faster with minimal data.

\noindent\circnum{2}~\textbf{Contrastive Learning.} ECDP~\cite{yang2023ecdp} (ICCV'23) used paired event-RGB data with conditional masking and probability distribution alignment for pre-training. Its extension, ECDP-Dense~\cite{yang2024ecdpdense} (ECCV'24), generalized to dense prediction tasks (segmentation, flow, depth) by enforcing context-to-context similarity. CM3AE~\cite{wang2025cm3ae} (ACM MM 2025) built a unified pre-training framework across RGB frames, event images, and event voxels using 2.5M data pairs with multi-modal fusion reconstruction and contrastive learning.

\noindent\circnum{3}~\textbf{Temporal Modeling.} TESPEC~\cite{liu2025tespec} (ICCV'25) introduced temporal tube masking across event histograms and used a recurrent backbone to reconstruct intensity video at masked patches --- the first method to explicitly leverage long event sequences during pre-training. STP~\cite{yang2025stp} (ICCV'25) achieved remarkable efficiency: 68.87\% top-1 accuracy on N-ImageNet with 1/10 of the pre-training parameters and 1/3 of the training epochs compared to prior methods.

\noindent\circnum{4}~\textbf{Physics-Inspired and Task-Agnostic.} EventPretrain~\cite{wang2025eventpretrain} (ACM MM 2025) proposed a three-stage physics-inspired framework: Difference-guided Masked Modeling, Backbone-fixed Feature Transition, and Focus-aimed Contrastive Learning, showing strong robustness to sparse and noisy events. EvRepSL~\cite{kang2024evrepsl} trained a representation generator that, once learned, produces device- and task-agnostic representations improving downstream performance without fine-tuning. SD2Event~\cite{huang2024sd2event} (CVPR'24) learned self-supervised dynamic keypoint detectors and contextual feature descriptors.

\noindent\circnum{5}~\textbf{Cross-Modal Knowledge Transfer.} EventDance~\cite{zheng2024eventdance} (CVPR'24) was the first to achieve source-free image-to-event cross-modal adaptation --- transferring knowledge from an image-trained model to events without accessing the source image data. EventFly~\cite{kong2025eventfly} (CVPR'25) addressed cross-platform adaptation, enabling a model trained on vehicle-mounted events to transfer to drones and quadruped robots through Event Activation Prior (EAP), EventBlend data mixing, and EventMatch dual-discriminator alignment. Depth AnyEvent~\cite{xu2025depthanycvent} (ICCV'25) distills dense proxy depth labels from Depth Anything v2 to train event-specific depth networks, achieving competitive results with fully supervised methods \textit{without requiring any depth annotations}. S5-ViT~\cite{zubic2024s5vit} (CVPR'24 Spotlight) demonstrated that SSM architectures pre-trained on events transfer robustly across different input frequencies, addressing a key practical challenge in event deployment. FFEvent~\cite{song2026ffevent} introduced Fourier-based RGB-to-event transfer with a bidirectional reverse state space model, achieving state-of-the-art across 8 benchmarks.

\textit{Discussion.} The progression from closed-set perception (\cref{sec:perception}) to open-vocabulary understanding represents a paradigm shift: instead of collecting labeled event data for each new category, a single model trained with language supervision can recognize any category. This dramatically reduces the annotation burden and enables deployment in open-world scenarios where the set of target objects is not known in advance. The rapid progress in event SSL raises an important question: should the field pursue \textit{event-native} foundation models trained from scratch on event data, or continue adapting RGB-pre-trained models? Current evidence suggests that adaptation is more data-efficient, but native pre-training --- if sufficient event data becomes available --- may ultimately capture event-specific temporal structure that RGB priors cannot provide.


% ----------------------------------------------------------------------
% 3.3.2 Scene Understanding and Reasoning
% ----------------------------------------------------------------------
\subsubsection{Scene Understanding and Reasoning}
\label{sec:understanding_reasoning}

While open-vocabulary perception (\cref{sec:understanding_ov}) endows event models with the ability to recognize \textit{what} is in a scene, scene understanding and reasoning address \textit{why} and \textit{how} --- answering questions, grounding language in event streams, and enabling embodied agents to act on event-based perception. We organize this subsection around multimodal LLMs, visual grounding and embodied intelligence, and language-guided recognition.

\paragraph{Event-Based Multimodal LLMs.}
Multimodal large language models (MLLMs) have demonstrated remarkable visual understanding capabilities on images and videos. Extending MLLMs to process event camera data opens the door to scene-level reasoning, visual question answering, and free-form description of dynamic event streams --- capabilities far beyond what task-specific models can provide.

EventGPT~\cite{zhou2025eventgpt} (CVPR'25) was the first MLLM designed for event stream understanding. It employs an event encoder, a spatio-temporal aggregator, and an event-language adapter, trained through a three-stage paradigm: (i)~warm-up on RGB-text pairs, (ii)~alignment on synthetic N-ImageNet-Chat data ($\sim$1M pairs), and (iii)~fine-tuning on real-world Event-Chat data. EventGPT supports scene summarization, reasoning, and question answering on event streams. EventVL~\cite{li2025eventvl} introduced an Event Spatiotemporal Representation for comprehensive feature extraction, together with EventVL-Base (1.4M multimodal pairs) and EventVL-QA (75K QA pairs) --- the largest event-language datasets to date. EventFlash~\cite{kong2026eventflash} addresses the efficiency bottleneck: current event MLLMs process only short event clips (\eg, 5 temporal bins in EventGPT), whereas EventFlash supports up to 1,000 bins through adaptive temporal window aggregation and sparse density-guided attention, achieving a 12.4$\times$ throughput improvement. Its EventMind dataset contains over 500K instruction sets.

An alternative direction explores whether existing LLMs can understand events \textit{without any event-specific training}. Wu~\etal~\cite{wu2025llmzeroshot} (ICASSP'25) showed that GPT-4o, when combined with CLIP-based event-to-image conversion and carefully designed prompts, achieves zero-shot event recognition that surpasses supervised baselines on N-ImageNet. LLM-EvRep~\cite{cho2025llmevrep} (WWW'25) trains an encoder-decoder to generate LLM-compatible event representations via self-supervised learning, outperforming E2VID-based conversion by 15--50\% when evaluated with GPT-4o.

\textit{Discussion.} Event MLLMs are still in their infancy. The largest event-language dataset (EventMind, 500K+) is orders of magnitude smaller than image-language resources (LAION-5B). Temporal reasoning over long, continuous event streams remains largely unexplored. Yet the early results demonstrate that MLLMs can bring a qualitatively different level of understanding to event data --- moving from ``what object?'' to ``what is happening, why, and what might happen next?''

\paragraph{Visual Grounding and Embodied Intelligence.}
Language grounding --- connecting event streams with natural language descriptions --- enables humans to query event-based systems using free-form text, while embodied intelligence applies these capabilities to real-world agents.

Talk2Event~\cite{zheng2025talk2event} (NeurIPS'25) established the first large-scale benchmark for language-driven object grounding in event-based perception, containing 5,567 scenes, 13,458 annotated objects, and over 30,000 validated referring expressions with four structured grounding attributes (appearance, status, relation-to-viewer, relation-to-others). It proposes EventRefer with a Mixture of Event-Attribute Experts (MoEE) architecture. NeuroViG~\cite{li2025neurovig} (WACV'25) used event streams for selective triggering of RGB cameras in video grounding, achieving 4$\times$ energy reduction with less than 1\% accuracy loss.

On the frontier of embodied intelligence, EP-VLM~\cite{yang2025epvlm} uses event-camera motion priors to guide patch-wise sparsification of RGB visual inputs for VLMs, achieving 50\% FLOPs savings while retaining 98\% accuracy. E-VLA~\cite{miao2026evla} introduces an event-augmented Vision-Language-Action model for robotic manipulation under extreme conditions, maintaining over 80\% task success in complete darkness where image-only baselines fail entirely --- demonstrating the practical value of events for embodied AI.

\paragraph{Language-Guided Recognition.}
Ev-LaFOR~\cite{cho2023evlafor} (ICCV'23 Oral) pioneered label-free event recognition by reconstructing images from events and performing recognition through CLIP, using category-guided attraction and repulsion losses to bridge textual and visual features. EventDance++~\cite{zheng2024eventdancepp} extended cross-modal adaptation with language-guided VLM enrichment for source-free event recognition.

\cref{tab:understanding_catalog} provides a comprehensive catalog of all scene understanding and reasoning methods.


% ----------------------------------------------------------------------
% 3.3.3 Event Data Simulation and Generation
% ----------------------------------------------------------------------
\subsubsection{Event Data Simulation and Generation}
\label{sec:understanding_gen}

A critical bottleneck for event camera research is the scarcity of large-scale, diverse, and accurately labeled event data. This subsection covers three complementary approaches to addressing this bottleneck: physics-based simulators that generate synthetic events from existing video or 3D environments, learned generative models that produce events from text or other conditional inputs, and sim-to-real transfer methods that bridge the gap between synthetic and real event data.

\paragraph{Physics-Based Simulators.}
Physics-based event simulators convert conventional video or rendered frames into synthetic event streams by modeling the logarithmic intensity change and temporal dynamics of event camera pixels. ESIM~\cite{rebecq2018esim} established the foundational framework for event camera simulation, generating events from an adaptive rendering engine that interpolates camera trajectories and computes per-pixel brightness changes. v2e~\cite{hu2021v2e} extended simulation to arbitrary video inputs with realistic noise modeling, enabling the conversion of any existing video dataset into synthetic event data. DVS-Voltmeter~\cite{lin2022dvsvoltmeter} proposed a more physically accurate model by incorporating pixel-level circuit dynamics including bandwidth limitations and refractory periods. These simulators have been instrumental in generating training data for supervised learning (\eg, the training of E2VID~\cite{rebecq2019e2vid} and many subsequent methods), creating evaluation benchmarks with ground-truth labels, and enabling large-scale self-supervised pre-training.

\paragraph{Learned and Neural Generation.}
A creative application of foundation models is generating synthetic event data from text or other conditional inputs, addressing the data scarcity problem from the generative side.

Text-to-Events~\cite{barchid2024texttoevents} (NICE'24) was the first text-to-events model, generating event streams of smooth motion directly from text prompts without intermediate intensity frame generation, using an autoencoder plus diffusion model architecture. ControlEvents~\cite{wu2025controlevents} leveraged Stable Diffusion priors for controllable event data synthesis conditioned on text labels, 2D skeletons, and 3D body poses. Notably, it can generate events from \textit{unseen text labels}, enabling a form of zero-shot data augmentation.

\textit{Discussion.} Text-guided event generation is still nascent but holds significant potential. If synthetic events can faithfully replicate the statistical properties of real events, they could dramatically scale up training data for all event-based tasks --- particularly for rare or dangerous scenarios (\eg, near-collisions, extreme weather) that are difficult to capture with real event cameras.

\paragraph{Sim-to-Real Transfer.}
Despite the availability of simulators, a significant domain gap persists between synthetic and real events due to noise characteristics, contrast threshold variations, and temporal dynamics that are difficult to model accurately. Bridging this sim-to-real gap is essential for methods trained on simulated data to generalize to real-world deployment. E2VID+~\cite{stoffregen2020hqf} addressed this gap by training on more realistic simulated data and introducing a principled photometric loss. EventFly~\cite{kong2025eventfly} (CVPR'25) introduced cross-platform adaptation strategies including EventBlend data mixing, enabling transfer across different event camera platforms and mounting configurations. More broadly, the self-supervised and cross-modal transfer methods discussed in \cref{sec:understanding_ov} provide indirect sim-to-real bridges by learning representations that are robust to domain shifts.

\cref{tab:understanding_catalog} provides a comprehensive catalog of all event data simulation and generation methods.
